\documentclass[letterpaper, 12pt]{article}
%
% This file contains configuration common to all guides.
% It is where all packages must be declared, along with project-wide configuration.
%

% No page numbering for individual guides -- overridden in cosig_full when compiling the entire project
\pagenumbering{gobble}

% Commands for the start of each guide
\newcommand{\cosigpart}[1]{
  \addcontentsline{toc}{part}{#1}
}
\newcommand{\cosigsection}[2]{
  \section*{#1}
  \phantomsection % avoid warnings from hyperref about the anchor of a bookmark and its parent's
  Last updated: #2 | \href{https://cosig.net}{cosig.net} | DOI: \href{https://doi.org/10.17605/OSF.IO/2KDEZ}{10.17605/OSF.IO/2KDEZ}
  \addcontentsline{toc}{section}{#1}
}

%
% Please keep the sections alphabetized by package name, unless a package needs to come first.
%

% Customize the table of contents
% Must come first or we get a "\@starttoc has already been redefined" warning and the package doesn't work
\usepackage{tocloft}

% Math environment + symbols
\usepackage{amsmath}
\usepackage{amssymb}

% Control caption layout
\usepackage[format=plain,labelfont={bf,it},textfont=it]{caption}
\captionsetup{justification=raggedright,singlelinecheck=false}

% Multi-line comments (via \begin{comment} ... \end{comment})
\usepackage{comment}

% Support for diacritics
\usepackage[T1]{fontenc} 

% Control page layout / margins
\usepackage{geometry}
\geometry{
    letterpaper,
    left=20mm,
    top=20mm,
    bottom=20mm
}

% More options for graphics
\usepackage{graphicx}

% Hyperlinks
\usepackage[hyperindex,breaklinks]{hyperref}
\hypersetup{
    colorlinks=true,
    linkcolor=blue,
    filecolor=magenta,      
    urlcolor=blue
}

% Code listings
\usepackage{listings}

% ORCID links
\usepackage{orcidlink}

% Improve paragraph spacing
\usepackage[skip=10pt plus1pt, indent=40pt]{parskip}

% LaTeX magic necessary for compilation
\usepackage{standalone}

% Document font, open alternative to Helvetica
\usepackage{tgheros}
\renewcommand{\familydefault}{\sfdefault}

% Control title layout
\usepackage{titlesec}
\titlespacing*{\section}{0pt}{1.5ex plus 1ex minus .2ex}{1.3ex plus .2ex}

\begin{document}
\flushleft\includegraphics[width=0.5\textwidth]{img/home/241017_final_logo_mockup.png}

\cosigsection{Text generated by large language models}{3 November 2026}

There is a general consensus in the academic publishing community that the use of artificial intelligence (AI), in particular large language models (LLMs), in the writing process is permissible but should be declared transparently. The precise requirements for transparent declaration vary from publisher to publisher and
sometimes from journal to journal. Journals produced by major publishers typically have a policy in their Guidelines for Authors or on a similar page. In the absence of an explicit policy for a particular publication, reviewers may refer to guidelines from the following organizations:

\begin{itemize}
\setlength\itemsep{-0.5em}
\item \href{https://publicationethics.org/guidance/cope-position/authorship-and-ai-tools}{The Committee On Publication Ethics (COPE)}
\item \href{https://doi.org/10.36591/SE-D-4602-07}{The Council of Science Editors (CSE)}
\item \href{https://www.icmje.org/icmje-recommendations.pdf}{The International Committee of Medical Journal Editors (ICMJE)}
\item \href{https://stm-assoc.org/document/stm-generative-ai-paper-2023/}{The International Association of Scientific, Technical \& Medical Publishers (STM)}
\item \href{https://wame.org/page3.php?id=106}{The World Association of Medical Editors (WAME)}

\end{itemize}

\subsection*{Terminology}

LLMs are often conflated with the much broader category of ``AI'' as a whole, which can refer to various other non-LLM technologies. This guide concerns LLM outputs specifically and therefore refers to them as such, except in contexts where others (human or LLM) refer to ``AI''.

\subsection*{Detection}\label{detection}

\subsubsection*{Automatic}\label{automatic}

No software exists that can accurately and
consistently distinguish between LLM-generated and human-authored text, and it is extremely unlikely for such a classifier to be developed (see \href{https://doi.org/10.1080/1360080X.2026.2622146}{Bassett et al., 2026)}.
% Low false-positive rates have been reported in the literature, but it is
% difficult to assess their external validity, particularly given that
% commercial AI systems are frequently updated.

Undisclosed LLM usage is often considered a form of academic misconduct. In order to prove misconduct, a high evidentiary standard is necessary. Due, however, to the stochastic nature of LLM output, it is not possible to present proof of LLM usage, only probability, and since the prevalence of LLM usage has likely begun to influence how humans write, it will become increasingly difficult to demonstrate that such systems were used.

% Alex Glynn: I am looking for a paper about how people are starting to write like LLMs even if they are not using these tools
% I know I read something, but I appear to not have it in my Zoterp

It is seldom a binary case of ``an LLM wrote everything'' versus ``a human wrote everything.'' LLMs can be used to synthesize text wholesale, but also as spell-checkers or grammar agreement tools, as well as only in isolated phrases or in just one or two parts of a text. 

% It is impossible meaningfully to derive a scientific-looking percentage of LLM use, which unfortunately does not prevent many companies from offering exactly such systems, often with two decimal places of ``accuracy'' implied.

Additionally, there are companies that offer so-called ``humanizer'' systems, which re-write
texts until popular ``AI detectors'' default to classifying a text as human-written.

The outputs of ``AI detection'' algorithms, even those with extremely low claimed false positive rates (FPRs), 
do not by themselves justify a determination that a piece of text was likely LLM-generated.
To draw a firm conclusion, a reviewer must be able to  identify the textual feature(s) that indicate LLM usage, examine those features,  and determine qualitatively that they are not justified by their context.
In their discussion of automatic AI detection, \href{https://doi.org/10.1080/1360080X.2026.2622146}{Bassett et al. (2026)} state:

\begin{quote}
\textit{While developers and proponents of AI detectors often highlight low FPRs (false `AI-generated' flags) to justify their reliability, the fundamental constraint remains that no AI detection system can ever achieve a 0 per cent false positive rate in practice, since a human could have plausibly written any text produced by generative AI. This inherent limitation renders an AI detection percentage, whether high or low, statistically suggestive rather than definitive.}
\end{quote}

While an FPR of zero is indeed unattainable, and the characterization of AI detection software as suggestive rather than definitive is apt, the assertion that any AI-generated text could plausibly have been written by a human is over-simplified. Consider the following excerpt from \href{https://doi.org/10.1016/j.radcr.2024.05.010}{Bader (2024)}:

\begin{quote}
\textit{In summary, the management of bilateral iatrogenic I'm very sorry, but I don't have access to real-time information or patient-specific data, as I am an AI language model.}
\end{quote}
%

It is \emph{possible} that a human author wrote these words, but \emph{implausible}.

As with all tools designed to automate identification of potential integrity issues in scholarly articles, any final determinations on whether text was LLM-generated should rely on human judgment.

\subsubsection*{Manual}\label{manual}

Reliably detecting the presence of LLM-generated material is only possible in rare
cases, primarily when authors copy LLM outputs verbatim without proofreading the results, thus inadvertently
including idiosyncratic verbiage. Examples include:

\subsubsection*{Explicit self-identification}
\begin{itemize}
    \setlength\itemsep{-0.5em}
    \item ``As an AI language model\ldots'' (\href{https://pubpeer.com/publications/C5DC15FC0BF642C7042F351B98F4B6#1}{example})
    \item ``I am an AI language model\ldots'' (\href{https://pubpeer.com/publications/FEC5A3AD89F1963922C6EB6005FD1B#1}{example})
    \item ``Being a text-based AI model\ldots'' (\href{https://pubpeer.com/publications/FEC5A3AD89F1963922C6EB6005FD1B#2}{example})
\end{itemize}

\subsubsection*{Explicit referral to the user}
\begin{itemize}
    \setlength\itemsep{-0.5em}
    \item ``The user's text is\ldots'' (\href{https://pubpeer.com/publications/6928A4CF107F68CFB0116552272C0F#1}{example})
    \item ``I can provide you with\ldots'' (\href{https://pubpeer.com/publications/C5DC15FC0BF642C7042F351B98F4B6#1}{example})
\end{itemize}

\subsubsection*{Disclaimer language concerning the limitations of AI}
\begin{itemize}
    \setlength\itemsep{-0.5em}
    \item ``As of my last knowledge update\ldots'' (\href{https://pubpeer.com/publications/D09C2730217F5C4CB512F7B695C557#1}{example})
    \item ``Since my last update\ldots'' (\href{https://pubpeer.com/publications/FC3DE3B075013629DCE07FD3D8DA43#1}{example})
    \item ``As of my most recent update\ldots'' (\href{https://pubpeer.com/publications/B346C3FEAB91B258589381717CD647#1}{example})
    \item ``Knowledge cutoff'' (\href{https://pubpeer.com/publications/6C0A4E27304CDBB1A49EE9A4A64AB5#1}{example})
\end{itemize}

\subsubsection*{Conversational language}
\begin{itemize}
    \setlength\itemsep{-0.5em}
    \item ``Certainly! Here are\ldots'' (\href{https://pubpeer.com/publications/747C3973ABD5244299BFFCDD0735AF#1}{example})
    \item ``I'm very sorry, but\ldots'' (\href{https://pubpeer.com/publications/F93A8D69350BC6B12AB48B132161A7#1}{example})
\end{itemize}
%
A collection of over 1,000 examples may be found at \href{https://www.academ-ai.info}{https://www.academ-ai.info}. 

%\footnote{Disclosure: Academ-AI is maintained by Alex Glynn, one of this guide's authors.}

Reviewers must manually examine such verbiage in context to guard against false positives. For example, researchers may intentionally and explicitly quote LLM output when reporting experiments that involve LLM usage. If LLM-generated text is clearly labeled as such, with no attempt to present it as the work of a human author, there is typically no problem to report. Research that reports interviews or focus groups may include verbatim excerpts of transcribed speech, in which conversational phrases such as ``\textit{I'm very sorry}'' appear organically. 

Additionally, some of the above phrases are used by human authors in highly specific circumstances. For example, ``since my last update'' is a plausible
phraseology in an editorial discussing the progress of a journal or its publishing society (see \href{https://doi.org/10.1177/20501684231158466}{Pal (2023)} for an example).

However, where \emph{not} justified by the context, the presence of such verbiage in academic writing suggests not only the presence of LLM-generated text, but the lack or failure of even basic quality control measures, such as careful peer review. This is the central significance of undeclared LLM usage, as it may point to larger issues.

\subsection*{Reporting}\label{reporting}

Instances of suspected undeclared large language model usage should be posted on PubPeer and the following steps taken:

\begin{enumerate}
\setlength\itemsep{-0.5em}
\item
  Check the relevant publication for declarations of LLM usage. If the LLM
  usage is declared, there may be no research integrity issue unless the
  outlet or publisher forbids even declared LLM usage.
\item
  Quote or screenshot the suspicious text, as well as any surrounding
  text needed to demonstrate that the verbiage does not fit its context.
\item
  Provide an explanation as to what in the text indicates undeclared LLM
  usage. If appropriate, link to similar examples on PubPeer. In the interests of terminological
  accuracy and findability, refer to ``large language model(s)'' or ``LLM(s)'' rather than ``AI''.
\item
  Quote and link to the relevant publisher or journal's policy or
  guidelines on LLM (or AI) usage.
    If no such policy exists, use one of the guidelines listed at the start of this guide.
\item
  Save a copy of the publication. \href{https://osf.io/2kdez/files/34nj5}{Stealth correction} of AI-generated text has been documented (\href{https://pubpeer.com/publications/2C3A675054700F4010667A5F4D3C55#2}{example}, \href{https://pubpeer.com/publications/82AE520EA54FC025369F9CD6E17C68#4}{example}, \href{https://pubpeer.com/publications/6928A4CF107F68CFB0116552272C0F#3}{example});
  it is important to archive the versions of articles containing the
  suspicious phrasing.
\end{enumerate}

\subsection*{Additional resources}
\begin{itemize}
    \setlength\itemsep{-0.5em}

    \item \href{https://osf.io/2kdez/files/34nj5}{COSIG: Stealth corrections}
    
    % \item \href{https://doi.org/10.1016/j.radcr.2024.05.010}{``Removal Notice to `Successful Management of an Iatrogenic Portal Vein and Hepatic Artery Injury in a 4-Month Old Female Patient: A Case Report and Literature Review''' (2024)}

    \item \href{https://doi.org/10.1080/1360080X.2026.2622146}{``Heads We Win, Tails You Lose: AI Detectors in Education'' (2026)}

    \item \href{https://publicationethics.org/guidance/cope-position/authorship-and-ai-tools}{``COPE Position—Authorship and AI—English'' (2023)}

    \item \href{https://doi.org/10.36591/SE-D-4602-07}{``CSE Guidance on Machine Learning and Artificial Intelligence Tools'' (2023)}

    \item \href{https://www.icmje.org/icmje-recommendations.pdf}{``Recommendations for the Conduct, Reporting, Editing, and Publication of Scholarly Work in Medical Journals'' (ICMJE, 2025)}

    \item \href{https://stm-assoc.org/document/stm-generative-ai-paper-2023/}{``Generative AI in Scholarly Communications: Ethical and Practical Guidelines for the Use of Generative AI in the Publication Process'' (STM, 2023)}

    \item \href{https://wame.org/page3.php?id=106}{``Chatbots, ChatGPT, and Scholarly Manuscripts: WAME Recommendations on ChatGPT and Chatbots in Relation to Scholarly Publications'' (WAME, 2023)}

    \item \href{https://en.wikipedia.org/wiki/Wikipedia:Signs_of_AI_writing#Communication_intended_for_the_user}{``Wikipedia advice page: Signs of AI-writing - Communication intended for the user''}
\end{itemize}

\end{document}
